\section*{Hoofdstuk \ref{ch:g12:climate-change} --- Klimaatverandering: waarnemingen, oorzaken en projecties}

\begin{solution}{exo:g12:climate-change:1}
De verandering die ze veroorzaakt in de netto energiestroom die het klimaatsysteem
binnenkomt aan de top van de atmosfeer, voordat de oppervlaktetemperatuur reageert;
gemeten in \unit{W/m^2}. Positief: de planeet wint energie en warmt op;
negatief: ze verliest energie en koelt af.
\end{solution}

\begin{solution}{exo:g12:climate-change:2}
Het hoogst in mei, het laagst in september. In de lente en de zomer neemt de begroeiing
van het noordelijk halfrond, waar de meeste landplanten groeien, door fotosynthese meer
koolstofdioxide op dan planten en bodems afgeven; in de herfst
en de winter overheersen afbraak en ademhaling en stijgt de concentratie weer.
\end{solution}

\begin{solution}{exo:g12:climate-change:3}
\emph{Zeer waarschijnlijk}: een kans van 90 tot 100\,\%; \emph{waarschijnlijk}: 66 tot
100\,\%. $[0.95 \text{ tot } 1.20]$ is een zeer waarschijnlijk bereik: 90\,\% kans dat
de ware waarde tussen 0.95 en 1.20 ligt.
\end{solution}

\begin{solution}{exo:g12:climate-change:4}
91\,\%. Water zet uit wanneer het opwarmt: een warmere oceaan neemt meer volume in, dus
stijgt zijn oppervlak zelfs als er geen water wordt toegevoegd (smeltend landijs voegt daar
nog water aan toe).
\end{solution}

\begin{solution}{exo:g12:climate-change:5}
Waterdamp: warmere lucht bevat meer waterdamp, een broeikasgas, wat de opwarming
versterkt. Sneeuw en ijs: smelten legt donkerdere oppervlakken bloot, die
meer zonlicht absorberen. De negatieve reactie: een warmere planeet straalt meer infrarood
uit naar de ruimte, wat haar terug in evenwicht brengt.
\end{solution}

\begin{solution}{exo:g12:climate-change:6}
$427.35 - 315.98 = \qty{111.37}{ppm}$; $111.37/315.98 \approx 0.352$, een
toename van ongeveer 35\,\%. Over 66 jaar: $111.37/66 \approx
\qty{1.7}{ppm}$ per jaar gemiddeld.
\end{solution}

\begin{solution}{exo:g12:climate-change:7}
$F = 3.93 \times \ln(427.35/278)/\ln 2 = 3.93 \times 0.4300/0.6931 \approx
\qty{2.44}{W/m^2}$.
\end{solution}

\begin{solution}{exo:g12:climate-change:8}
$3.7 \times 30 = \qty{111}{mm}$, ongeveer \qty{11}{cm}. $3.7/1.3 \approx 2.8$:
bijna drie keer sneller.
\end{solution}

\begin{solution}{exo:g12:climate-change:9}
$\Delta T = 2.72/1.16 \approx \qty{2.3}{\degreeCelsius}$. Het klimaat is niet
in evenwicht: de oceaan, die ongeveer 90\,\% van de extra energie opneemt,
warmt langzaam op, dus heeft het oppervlak de \omterm{def:g12:climate-change:forcing}{forcering} nog niet ingehaald (het
systeem wint nog altijd ongeveer \qty{0.8}{W/m^2}).
\end{solution}

\begin{solution}{exo:g12:climate-change:10}
De Zon en vulkanen verklaren hoogstens \qty{0.1}{\degreeCelsius} van een
opwarming van \qty{1.06}{\degreeCelsius}: ze kunnen er niet de oorzaak van zijn. De andere
natuurlijke oorzaak is de interne variabiliteit van het klimaat (zijn natuurlijke
schommelingen, zoals warmte-uitwisselingen tussen oceaan en atmosfeer); de
evaluatie schat haar bijdrage op $-0.2$ tot $+\qty{0.2}{\degreeCelsius}$,
ook veel te klein.
\end{solution}

\begin{solution}{exo:g12:climate-change:11}
Zeer waarschijnlijk onder \qty{2}{\degreeCelsius}: alleen SSP1-1.9 (zeer waarschijnlijk bereik
1.0 tot 1.8). Zeer waarschijnlijk erboven: SSP2-4.5 (2.1 tot 3.5), SSP3-7.0 (2.8 tot 4.6)
en SSP5-8.5 (3.3 tot 5.7). Voor SSP1-2.6 (1.3 tot 2.4) is geen van beide zeer waarschijnlijk.
\end{solution}

\begin{solution}{exo:g12:climate-change:12}
De aerosolforcering is $-0.22 - 0.84 = \qty{-1.06}{W/m^2}$; zonder haar wordt het
totaal $2.72 + 1.06 = \qty{3.78}{W/m^2}$. De evenwichtsopwarming
neemt toe met $1.06/1.16 \approx \qty{0.9}{\degreeCelsius}$.
\end{solution}

\begin{solution}{exo:g12:climate-change:13}
Omdat $\ln(2C) - \ln C = \ln 2$ voor elke $C$, geldt
\[
  F(556) - F(278) = \frac{3.93\,\ln 2}{\ln 2} = \qty{3.93}{W/m^2}, \qquad
  F(1112) - F(556) = \frac{3.93\,\ln 2}{\ln 2} = \qty{3.93}{W/m^2}:
\]
dezelfde \omterm{def:g12:climate-change:forcing}{forcering} voor elke verdubbeling. Totaal bij \qty{1112}{ppm}: $2 \times 3.93 =
\qty{7.86}{W/m^2}$.
\end{solution}

\begin{solution}{exo:g12:climate-change:14}
$\dfrac{\dd F}{\dd C} = \dfrac{F_{2\times}}{C \ln 2}$. Bij \qty{280}{ppm}:
$3.93/(280 \times 0.6931) \approx \qty{0.0202}{W/m^2}$ per ppm; bij
\qty{420}{ppm}: $\approx \qty{0.0135}{W/m^2}$ per ppm. Elke extra ppm voegt
minder \omterm{def:g12:climate-change:forcing}{forcering} toe naarmate de concentratie groeit (1.5 keer minder bij 420 dan bij 280),
maar de concentratie blijft stijgen, dus groeit de \omterm{def:g12:climate-change:forcing}{forcering} nog altijd.
\end{solution}

\begin{solution}{exo:g12:climate-change:15}
$3.93/3.22 \approx \qty{1.2}{\degreeCelsius}$ zonder \omterm{def:g12:climate-change:feedback}{terugkoppelingen};
$3.93/1.16 \approx \qty{3.4}{\degreeCelsius}$ ermee. De \omterm{def:g12:climate-change:feedback}{terugkoppelingen} versterken de opwarming
met een factor $3.22/1.16 \approx 2.8$.
\end{solution}

\begin{solution}{pb:g12:climate-change:1}
\textbf{1.} $3.37 + 3.51 = \qty{6.88}{ppm}$, ongeveer \qty{7}{ppm}.

\textbf{2.} Van de lente tot de late zomer halen de planten van het noordelijk
halfrond door fotosynthese koolstofdioxide sneller uit de lucht dan
ademhaling en afbraak het teruggeven: de concentratie daalt tot september.
In de herfst en de winter keert de balans om, en klimt de concentratie
terug naar haar maximum van mei.

\textbf{3.} Het zuidelijk halfrond heeft veel minder land, en dus veel minder
begroeiing, dan het noordelijk; het grootste deel van zijn oppervlak is oceaan.

\textbf{4.} Hoog gelegen ligt het boven de laag vervuilde lucht dicht bij de
grond; ver van continenten en steden is de lucht die het bemonstert goed gemengd en
vertegenwoordigt ze de hele atmosfeer, niet een plaatselijke bron.

\textbf{5.} $(325.68 - 316.91)/10 \approx \qty{0.88}{ppm}$ per jaar.

\textbf{6.} $(427.35 - 401.01)/10 \approx \qty{2.63}{ppm}$ per jaar.

\textbf{7.} $2.634/0.877 \approx 3.0$: de groeisnelheid is verdrievoudigd.

\textbf{8.} $(427.35 - 278)/278 \approx 0.537$: ongeveer 54\,\%.

\textbf{9.} $427.35 - 300 \approx \qty{127}{ppm}$.

\textbf{10.} $F = 3.93 \times \ln(427.35/278)/\ln 2 \approx
\qty{2.44}{W/m^2}$.

\textbf{11.} $F(316.91) = 3.93 \times \ln(316.91/278)/\ln 2 \approx
\qty{0.74}{W/m^2}$; toename $2.44 - 0.74 \approx \qty{1.70}{W/m^2}$.

\textbf{12.} Eén \unit{W/m^2} meer vermenigvuldigt de concentratie met
$2^{1/3.93} \approx 1.193$: $427.35 \times 1.193 \approx \qty{510}{ppm}$.

\textbf{13.} $(509.8 - 427.35)/2.634 \approx 31$ jaar.

\textbf{14.} $2 \times 278 = \qty{556}{ppm}$; \omterm{def:g12:climate-change:forcing}{forcering} $F_{2\times} =
\qty{3.93}{W/m^2}$.

\textbf{15.} Per definitie van de gevoeligheid \qty{3}{\degreeCelsius}
(waarschijnlijk 2.5 tot \qty{4}{\degreeCelsius}).

\textbf{16.} $\lambda = 3.93/3 \approx \qty{1.31}{W.m^{-2}.\degreeCelsius^{-1}}$.
De som van de \omterm{def:g12:climate-change:feedback}{terugkoppelingen}, \qty{1.16}{W.m^{-2}.\degreeCelsius^{-1}}, geeft
$3.93/1.16 \approx \qty{3.4}{\degreeCelsius}$, iets meer; de geëvalueerde
beste schatting combineert deze fysica met de waargenomen opwarming en klimaten van het
verleden, en beide waarden liggen in het waarschijnlijke bereik.

\textbf{17.} De oceaan slaat het grootste deel van de extra energie op en warmt langzaam op: het
oppervlak bereikt zijn evenwichtsopwarming pas na decennia tot eeuwen.

\textbf{18.} Ze neemt één verhaal over de toekomstige uitstoot aan (de groei van de
laatste tien jaar voor altijd aangehouden): dat is een \omterm{def:g12:climate-change:scenario}{scenario}. Een \omterm{def:g12:climate-change:scenario}{projectie} berekent de
reactie van het klimaat en van de koolstofkringloop op een \omterm{def:g12:climate-change:scenario}{uitstootscenario}
met fysische modellen, en geeft bereiken en waarschijnlijkheden.

\textbf{19.} $(556 - 427.35)/2.634 \approx 48.8$ jaar na 2025: ongeveer
\textbf{2074}.
\end{solution}
